\section{Appendix}
\subsection{List of Criteria}

The following criteria list $L = \{L_1, \dots, L_n\}$, originally
specified by Infosys, was used as the task specification governing the
content and output format for the summarisation systems (paraphrased here only for readability).

\begin{enumerate}[label=$L_{\arabic*}$:, leftmargin=2.2em]

\item \textbf{Role framing.} The model is instructed to act as an expert in
COBOL programming and documentation, tasked with analysing and categorising
a given COBOL code snippet by rule classification and rule type.

\item \textbf{No omission of code.} Every uncommented line of the input
must appear in the \texttt{Segregated Code} field, in the exact order of
the original source; no executable line may be skipped, even if it appears
redundant or deeply nested.

\item \textbf{Comment handling.} Lines beginning with comment indicators
(e.g.\ \texttt{*}, \texttt{*---}) or inline comments, and any block visually
separated as a comment, must be omitted from the segregated code.

\item \textbf{Logical segregation.} Code must be grouped by logical
function or condition (e.g.\ initialisation, adjustments, validations,
aging rules), preserving the original COBOL structure and nesting within
each group.

\item \textbf{Business rule description.} Each code block must be
accompanied by a business-readable description that explicitly covers
every line of the corresponding segregated code, in presentation order,
using a structured bullet hierarchy to
represent nested logic.

\item \textbf{Variable expansion.} All COBOL variable names must be
expanded into plain English, followed by the original COBOL identifier in
parentheses (e.g.\ ``total credit amount (TOT-CR-AMT-WS)'').

\item \textbf{Conditional coverage.} Every conditional path (\texttt{IF},
\texttt{ELSE}, \texttt{EVALUATE}, etc.) and its implications must be
explicitly included in the business rule explanation.

\item \textbf{Dependency-aware summarisation.} For every \texttt{PERFORM}
statement in the snippet, the model must incorporate the supplied business
summary of the corresponding dependent paragraph when summarising the
calling paragraph; if no \texttt{PERFORM} statements are present, no
statement about dependencies should be produced.

\item \textbf{Strict output schema.} Output must be a single JSON array of
objects, each containing exactly the fields: \texttt{Rule Type},
\texttt{Rule Classification}, \texttt{Paragraph Name}, \texttt{Segregated
Code} (array of strings), \texttt{Business Rule Extraction} (string), and
\texttt{Business Rule Extraction Summary} (string).

\item \textbf{Single-line string encoding.} All multi-line content
(e.g.\ business rule explanations) must be encoded as single-line JSON
strings using \texttt{\textbackslash n} as the line separator; no raw
multi-line text blocks are permitted.

\item \textbf{No empty or null fields.} Every field in every object must
be populated; empty strings or null values are disallowed.

\item \textbf{Output purity.} The response must contain only the JSON
array, no markdown, code fences, explanatory text, or extraneous
whitespace before or after the JSON structure.

\item \textbf{1:1 traceability.} Every line of code must map to exactly
one logical segment, and every logical segment's code must be fully and
uniquely accounted for in its accompanying business rule explanation.

\item \textbf{No interpretation beyond the source.} The model must not
alter, reinterpret, or infer missing context beyond what is present in
the COBOL code and the supplied dependency descriptions.

\end{enumerate}

%\subsection{Architecture and Prompts}